% Propietario conservado: /Users/ruben/Documents/New project/output/EXTREMA_MEDIA_RAZON_RECIPROCIDAD_ES_EN_20260918/ARTICULO/ES/source/nuclear/02.tex
\chapter{Conservación de observables y registro completo}
\label{x_reciprocidad:nuclear:2}
\section{Balance de un observable conservado}

Sean \(\mathcal H,\mathcal K\) las realizaciones de un espacio de lectura y de su extensión enriquecida. Sean \(J:\mathcal H\to\mathcal K\) una isometría, \(P=JJ^*\), \(U:\mathcal K\to\mathcal K\) un transporte acotado y \(\widetilde A\) un observable acotado autoadjunto. Supóngase

\[
U^*\widetilde A U=\widetilde A,\qquad [\widetilde A,P]=0.
\]

Definimos

\[
A=J^*\widetilde A J,\qquad T=J^*UJ,
\qquad \eta=(I-P)UJ.
\]

\begin{proposition}
\label{x_reciprocidad:nuc02:balance}
La carga se descompone exactamente entre lectura y complemento:

\begin{equation}
\boxed{A=T^*AT+\eta^*\widetilde A\eta.}\label{x_reciprocidad:nuc02:eq:1}
\end{equation}
\end{proposition}

\begin{proof}
La identidad \(UJ=JT+\eta\) es una descomposición ortogonal. Al expandir \(J^*U^*\widetilde A UJ\), los términos cruzados son cero porque \(\widetilde A\) reduce los subespacios \(P\mathcal K\) y \((I-P)\mathcal K\). El término izquierdo es \(A\); los dos términos diagonales son los de \eqref{x_reciprocidad:nuc02:eq:1}. Esto prueba también la identidad de las formas sesquilineales para dos estados distintos. Si \(\widetilde A\geq0\), ambos sumandos son positivos. Para \(\widetilde A=I\), la hipótesis de conservación es \(U^*U=I\).
\end{proof}

Si la lectura mezcla sectores del observable, el balance conserva adicionalmente

\[
T^*J^*\widetilde A\eta+\eta^*\widetilde A JT.
\]

Estos términos representan la interferencia entre las dos residencias de la carga. La condición de reducción escrita arriba garantiza su anulación para toda la descomposición. En realizaciones particulares pueden anularse también por otras razones, por ejemplo cuando \(\eta=0\).

\section{Especialización a la conexión nonádica}

Sea \(C\) la realización unitaria del transporte de memoria sobre \(\mathcal H\). En \(\mathcal H^9\), el transporte de las nueve fases y su inclusión uniforme son

\[
S_C(v_0,\ldots,v_8)=(Cv_8,v_0,\ldots,v_7),
\qquad Jv=\tfrac13(v,\ldots,v).
\]

El factor \(1/3\) resulta de la normalización de nueve copias. La compresión y su complemento son

\[
T=\frac{8I+C}{9},\qquad
\eta v=\frac1{27}\bigl(8(C-I)v,-(C-I)v,\ldots,-(C-I)v\bigr).
\]

Si \(A=A^*\) es acotado y \([A,C]=0\), el observable \(\widetilde A=I_9\otimes A\) satisface las hipótesis de la proposición~\ref{x_reciprocidad:nuc02:balance}. Por tanto,

\begin{equation}
\boxed{A=T^*AT+\frac8{81}(I-C)^*A(I-C).}\label{x_reciprocidad:nuc02:eq:2}
\end{equation}

Los coeficientes proceden del reparto entre las nueve componentes: \(8^2+8=72\), dividido por \(27^2\), da \(8/81\). La identidad es el balance sesquilineal ya documentado en el artículo espectral, expresado aquí para un observable general.

La iteración telescópica proporciona, para cada \(N\geq1\),

\begin{equation}
A=T^{*N}AT^N+
\sum_{j=0}^{N-1}T^{*j}\eta^*\widetilde A\eta T^j.\label{x_reciprocidad:nuc02:eq:3}
\end{equation}

Para \(A=I\), la aplicación

\[
V_Nv=(T^Nv,\eta v,\eta Tv,\ldots,\eta T^{N-1}v)
\]

es isométrica. El paso \(N\to N+1\) descompone exclusivamente el terminal \(T^Nv\) en su nueva lectura y complemento; mantiene todos los registros previos. Así, la conservación se verifica a cualquier profundidad.

Distinción de operaciones. Se cumple \(S_C^9=I_9\otimes C\). La iteración \(T^N\) describe compresiones sucesivas con almacenamiento de sus complementos. El transporte enriquecido sin compresiones intermedias es \(S_C^N\). Ambas operaciones tienen dominios definidos y funciones diferentes; la prueba conserva esa distinción.

\section{Sector terminal y registro a profundidad ilimitada}

Sea \(P_{\rm fix}\) la proyección ortogonal sobre \(\ker(I-C)\).

\begin{theorem}
Para todo transporte unitario \(C\), la compresión nonádica satisface

\begin{equation}
\underset{N\to\infty}{\operatorname{s-lim}}T^N=P_{\rm fix},\qquad
\boxed{I=P_{\rm fix}+\sum_{j=0}^{\infty}T^{*j}\eta^*\eta T^j.}\label{x_reciprocidad:nuc02:eq:4}
\end{equation}

La serie converge en la topología fuerte. Para cada observable acotado autoadjunto \(A\) que conmuta con \(C\), se obtiene además

\begin{equation}
\boxed{A=P_{\rm fix}AP_{\rm fix}
+\sum_{j=0}^{\infty}T^{*j}\eta^*(I_9\otimes A)\eta T^j.}\label{x_reciprocidad:nuc02:eq:5}
\end{equation}
\end{theorem}

\begin{proof}
En la representación espectral del operador unitario \(C\), la compresión corresponde a \(t(z)=(8+z)/9\). Sobre el círculo unitario,

\[
|t(z)|^2=1-\frac8{81}|1-z|^2.
\]

Por tanto \(t(z)^N\to0\) para \(z\ne1\), y \(t(1)^N=1\). La convergencia dominada respecto de cada medida espectral vectorial prueba \(T^N\to P_{\rm fix}\) fuertemente. El mismo argumento vale para \(T^{*N}\). Como \(A\) es acotado, \(T^{*N}AT^N\to P_{\rm fix}AP_{\rm fix}\) fuertemente. El límite de \eqref{x_reciprocidad:nuc02:eq:3} da \eqref{x_reciprocidad:nuc02:eq:5}, y su especialización \(A=I\) da \eqref{x_reciprocidad:nuc02:eq:4}.
\end{proof}

En particular,

\[
V_\infty v=(P_{\rm fix}v,\eta v,\eta Tv,\eta T^2v,\ldots)
\]

es una isometría. Conserva el producto interior completo, además de la norma.

\begin{corollary}[Memoria bilateral]
Para el desplazamiento \(Ce_n=e_{n+1}\) en \(\ell^2(\mathbb Z)\), \(P_{\rm fix}=0\): una sucesión fija sería constante y su única realización cuadrado-sumable es cero. En consecuencia, toda la norma inicial se distribuye en el registro de complementos. Esto precisa el comportamiento ilimitado de la realización bilateral de memoria presente en el corpus.
\end{corollary}

\begin{corollary}[Calendario finito]
En un ciclo de \(d\) posiciones, la sucesión uniforme determina un subespacio fijo unidimensional. Ese sector permanece como terminal y el complemento acumula el resto. El calendario finito y el registro bilateral poseen, por tanto, terminales diferentes.
\end{corollary}

En el caso bilateral, la convergencia fuerte puede coexistir con \(\|T^N\|=1\) para todo \(N\). La demostración no supone una tasa uniforme de contracción. También distingue esta compresión del modo \((5/9)C\), cuya contracción geométrica es uniforme y está estudiada separadamente en las fuentes.

\section{Conservación de simetría y transporte de carta}

Sea \(R\) unitario y definamos \(C'=RCR^*\), \(A'=RAR^*\). La fórmula explícita de la compresión da

\begin{equation}
T(C')R=RT(C),\qquad
\eta(C')R=R^{\oplus9}\eta(C).\label{x_reciprocidad:nuc02:eq:6}
\end{equation}

\begin{proof}
La primera igualdad se obtiene distribuyendo \((8I+RCR^*)R/9\). La segunda resulta de \((C'-I)R=R(C-I)\) en cada una de las nueve componentes de \(\eta\).
\end{proof}

La iteración de \eqref{x_reciprocidad:nuc02:eq:6} entrelaza los registros finitos. Además, \(P_{\rm fix}(C')=RP_{\rm fix}(C)R^*\), por transporte del núcleo de \(I-C\), de modo que también queda entrelazado el registro ilimitado. Así, el balance de observables es covariante cuando se transportan conjuntamente la carta, el operador y el observable. Las simetrías que conmutan con \(C\) actúan además en una misma carta fija. Esta distinción permite incorporar las transformaciones excepcionales con sus marcos, banderas y orientaciones efectivamente transportados.

La composición de \cref{x_reciprocidad:lem:k-hadamard,x_reciprocidad:exc:entrelazador,x_reciprocidad:exc:isometria} utiliza la matriz de Hadamard del registro, \(H_{12}^2=4I\), y la incidencia \(\mathcal I\) de las doce posiciones en las 132 hexadas:

\[
J_H=H_{12}/2,\qquad
\mathcal I^*\mathcal I=36I+30\mathbf1\mathbf1^*,\qquad
U_W=\mathcal I/6\big|_{\mathbf1^\perp}.
\]

Ambos operadores normalizados son isométricos en los dominios indicados. La composición \(B=U_WJ_H\) sobre \(W=J_H(\mathbf1^\perp)\) conserva productos interiores y entrelaza la representación transportada de \(M_{12}\) con la acción en las hexadas. El registro entero \(K\) conserva, además, la inversión integral \(H_{12}/4\) sobre su red de congruencias; esa inversión y la normalización isométrica tienen funciones distintas. La demostración, la prolongación APP–Witt–Leech y sus dominios están desarrollados aquí.

En esta composición, la conservación de incidencia tiene un contenido operatorio concreto. Su extensión por otros transportes usa \eqref{x_reciprocidad:nuc02:eq:6} y los entrelazadores efectivos de las fuentes, respetando sus grupos y dominios respectivos.

\subsection{Transporte íntegro de las doce coordenadas de K}

La separación del subespacio centrado puede completarse manteniendo también su media. Sean \(P_0=I-\mathbf1\mathbf1^*/12\), \(\mu(k)=\mathbf1^*k/12\) y

\[
F(k)=\left(\mu(k),\frac16\mathcal I P_0k\right)\in\mathbb R\oplus E_{\rm inc}.
\]

Equipamos el codominio con \(\|(\mu,y)\|^2=12|\mu|^2+\|y\|^2\). El Gram de incidencia demuestra

\begin{equation}
\|F(k)\|^2=\|k\|^2,\qquad
\boxed{F^{-1}(\mu,y)=\mu\mathbf1+\frac16P_0\mathcal I^*y.}\label{x_reciprocidad:nuc02:eq:6a}
\end{equation}

En efecto, \(\mathcal I^*\mathcal I P_0=36P_0\), de modo que la fórmula inversa recupera \(P_0k\), y el primer término recupera su media. La acción de \(M_{12}\) es trivial sobre la media y permuta las hexadas; por ello \(F\) es equivariante. Componer \(FJ_H\) transporta las coordenadas firmadas normalizadas. Así se conservan las doce coordenadas del registro, no sólo su parte centrada. El dominio de esta reconstrucción es el registro dodecafásico; la memoria completa de la historia mantiene sus otras coordenadas y su prolongación.

\subsection{Especialización al operador excepcional de orden tres}

La construcción excepcional produce el operador ortogonal \(g_c\) sobre la realización del vecino de Leech, mediante las doce fibras APP–Witt de tipo \(A_2\), y demuestra \(g_c^3=I\) y \(\ker(I-g_c)=0\). Su construcción, entrelazamiento y preservación reticular están localizados en \cref{x_reciprocidad:nuclear:fibras-app-witt,x_reciprocidad:nuc04:c3}.

Construimos ahora, como composición operatoria de esa salida con la carta de nueve fases,

\[
T_3=\frac{8I+g_c}{9},\qquad
\eta_3=(I-JJ^*)S_{g_c}J.
\]

\begin{corollary}
Esta composición satisface

\begin{equation}
\boxed{T_3^*T_3=\frac{19}{27}I,\qquad
\eta_3^*\eta_3=\frac8{27}I,\qquad
\|T_3^Nv\|^2=\left(\frac{19}{27}\right)^N\|v\|^2.}\label{x_reciprocidad:nuc02:eq:6b}
\end{equation}
\end{corollary}

\begin{proof}
El vector \((I+g_c+g_c^2)v\) es fijo por \(g_c\); la ausencia de vectores fijos implica \(I+g_c+g_c^2=0\). Como \(g_c^*=g_c^2\), tenemos \(g_c+g_c^*=-I\). Sustituyendo,

\[
T_3^*T_3=\frac{65I+8(g_c+g_c^*)}{81}
=\frac{57}{81}I=\frac{19}{27}I.
\]

El balance \eqref{x_reciprocidad:nuc02:eq:2} con \(A=I\) da el segundo coeficiente. La aplicación reiterada de la igualdad de normas da la potencia para todo \(N\).
\end{proof}

El registro ilimitado es isométrico y su terminal tiende a cero con la tasa explícita de \eqref{x_reciprocidad:nuc02:eq:6b}. Aquí la estructura excepcional determina los dos coeficientes complementarios; ambos proceden de la relación polinómica de \(g_c\) y de las nueve fases. Esta es una composición matemática desarrollada en esta investigación. Su definición no identifica \(g_c\) con \(\Gamma_9\), ni afirma que \(S_{g_c}\) sea el operador temporal canónico del TPK: conserva la distinción entre prolongación excepcional y holonomía temporal mientras las compone explícitamente.

\section{Refinamiento de historias y compatibilidad de los balances}

Para las historias enriquecidas de profundidad \(n\), sea \(p_h(\varepsilon)\) la probabilidad de una extensión superviviente, normalizada por \(\sum_\varepsilon p_h(\varepsilon)=1\). El refinamiento canónico sobre las bases de historias es

\[
R_n\delta_h=\sum_{\varepsilon\in\operatorname{Out}(h)}
\sqrt{p_h(\varepsilon)}\,\delta_{h\varepsilon}.
\]

Las extensiones de historias diferentes tienen prefijos diferentes y soportes disjuntos. Por ello \(R_n^*R_n=I\). Esta es la realización isométrica de la consistencia proyectiva de la medida: los coeficientes proceden de los pesos de supervivencia ya construidos.

Para transportar \eqref{x_reciprocidad:nuc02:eq:1} entre profundidades, se conservan conjuntamente las inclusiones, las proyecciones, los transportes y los observables. Con isometrías de refinamiento \(R_n,\widetilde R_n\), los cuadrados necesarios son

\[
J_{n+1}R_n=\widetilde R_nJ_n,\quad
P_{n+1}\widetilde R_n=\widetilde R_nP_n,\quad
U_{n+1}\widetilde R_n=\widetilde R_nU_n,\quad
\widetilde A_{n+1}\widetilde R_n=\widetilde R_n\widetilde A_n.
\]

Estos cuadrados implican \(T_{n+1}R_n=R_nT_n\) y \(\eta_{n+1}R_n=\widetilde R_n\eta_n\). Con cotas uniformes, las identidades pasan al límite inductivo de las realizaciones. Las historias, en cambio, forman el sistema proyectivo de prefijos; las direcciones contrarias de ambos tipos de mapas quedan diferenciadas. Se incluyen las pruebas de medida, Gram y paso al límite.

\section{Especialización variacional: acción y orientación}

La acción elíptica de \cref{x_reciprocidad:iv:ellipse:section,x_reciprocidad:nuclear:7} proporciona una especialización física del principio, con los parámetros construidos en sus etapas HMT. En sus coordenadas canónicas,

\[
H_\eta(Q,P)=\frac{\omega}{2}
\left(e^{-\eta}Q^2+e^\eta P^2\right),\qquad
I_\eta=H_\eta/\omega.
\]

La transformación canónica \(Q=e^{\eta/2}q\), \(P=e^{-\eta/2}p\) convierte \(I_\eta\) en \((q^2+p^2)/2\). Las rotaciones de \((q,p)\), transportadas a \((Q,P)\), constituyen una simetría continua. Para el lagrangiano \(L=P\dot Q-\omega I_\eta\), su carga de Noether es \(I_\eta\). En el contorno de la fuente \(I_\eta=\hbar\), la acción orientada es \(\oint P\,dQ=2\pi\hbar=h\), con la orientación de la órbita hamiltoniana. La prueba distingue la transformación simpléctica y la antisimpléctica: pueden conservar la misma energía y actuar con signos diferentes sobre la acción orientada. La demostración variacional completa se presenta en \cref{x_reciprocidad:nuclear:7}.

\subsection{Actualización de la forma y término de conexión}

La matriz de forma \(M_\eta=\operatorname{diag}(e^{\eta/2},e^{-\eta/2})\) permite desarrollar un transporte entre elipses con parámetros diferentes. Fijamos \(J_2=\begin{pmatrix}0&-1\\1&0\end{pmatrix}\) y \(R(\theta)=\exp(-\theta J_2)\), la rotación hamiltoniana en el plano \((q,p)\). El transporte

\[
z'=M_{\eta'}R(\delta\theta)M_\eta^{-1}z,
\qquad z=(Q,P)^\mathsf T,
\]

es simpléctico y satisface exactamente \(I_{\eta'}(z')=I_\eta(z)\). La prueba consiste en llevar ambos estados a sus coordenadas normalizadas y usar la conservación de \(q^2+p^2\) por la rotación. Conserva la acción mientras cambia la forma.

Para una realización diferenciable \(\eta(t)\), con avance angular \(\dot\theta=\omega\), la diferenciación de ese transporte produce

\[
\dot Q=\omega e^\eta P+\tfrac12\dot\eta Q,\qquad
\dot P=-\omega e^{-\eta}Q-\tfrac12\dot\eta P.
\]

Su generador hamiltoniano, obtenido de estas dos ecuaciones, es

\begin{equation}
\boxed{H_{\rm trans}=\omega I_\eta+\frac{\dot\eta}{2}QP.}\label{x_reciprocidad:nuc02:eq:7}
\end{equation}

El segundo sumando es el término de conexión de la forma variable. Con la convención \(\{f,g\}=\partial_Qf\partial_Pg-\partial_Pf\partial_Qg\),

\[
\partial_tI_\eta=\frac{\dot\eta}{2}(-e^{-\eta}Q^2+e^\eta P^2),
\quad
\left\{I_\eta,\frac{\dot\eta}{2}QP\right\}
=\frac{\dot\eta}{2}(e^{-\eta}Q^2-e^\eta P^2).
\]

Ambas contribuciones se cancelan y \(dI_\eta/dt=0\). Para el transporte prescrito, \eqref{x_reciprocidad:nuc02:eq:7} es único salvo una función del tiempo. La realización diferenciable es una especialización declarada del transporte entre formas; el resultado discreto anterior vale directamente entre dos parámetros generados. La actualización TPK específica conserva su propia regla para seleccionar esos parámetros y sus tiempos.

Así se obtiene una conexión precisa: un observable conservado por la dinámica completa se distribuye entre lectura y memoria mediante \eqref{x_reciprocidad:nuc02:eq:1}. Cuando ese observable es el generador de una simetría continua de la acción construida, la cantidad distribuida es una carga de Noether. La norma, la carga de una simetría y el funcional energético quedan relacionados por sus operadores específicos.

Una especialización acotada efectiva se obtiene del calendario de 108 estados. Su Hamiltoniano realizado es \(H_{\rm clk}=\hbar\omega_0N_{108}\), con \(N_{108}=\sum_{k=0}^{107}kP_k\) y \(P_k\) sus proyectores de Fourier. La acción del reloj tiene carga \(\hbar\langle\psi,N_{108}\psi\rangle=E_{\rm clk}/\omega_0\). Los transportes y las inclusiones que satisfacen las hipótesis de \eqref{x_reciprocidad:nuc02:eq:1} distribuyen esa carga entre lectura y complemento. Aquí todos los operadores son finitos y acotados. La acción elíptica de \eqref{x_reciprocidad:nuc02:eq:7} se ha probado en su realización clásica simpléctica; una cuantización con observables no acotados conserva, además, sus condiciones de dominio propias.
